

\section{Errors Clusters}


\ctt{ We should mention, that in any time that we call to $H^{\otimes n}$, even if its only for moving to the dual space, we immediately apply after the $\phi$-transpose mapping. The idea is that it will allow to us to ensure that the $X$-errors are always bounded by upper-right tingle and the $Z$-errors are always bounded by lower-left triangles.} 

\begin{definition}[$\Size{ }$-Function]
Consider a $(d,d^{\prime})$-Toric complex $\mathcal{G}$, and let $U$ be a subset of vertices, $U \subset \mathcal{G}$. We define the \textit{size} of $U$ to be:
  \begin{equation*}
    \begin{split}
      \Size{U} \colonequals \sum_{i}^{d+1} \max_{v \in S} L_{i}(v) 
    \end{split}
  \end{equation*}
Furthermore, we extend the notion to binary assignment over the $i$th faces of a complex $z : \mathcal{F}^{(i)} \rightarrow \mathbb{F}_{2}$ by associating $\Size{z}$ with the size of the vertex subset supporting $z$.
\end{definition}

  \begin{claim}
    Let $\mathbf{U} = \{ U_{1}, U_{2} , .. , U_{k} \}$ where $U_{i} \subset \mathbb{R}^{d}$, introduce an edge between $U_{i}$ and $U_{j}$ iff $U_{i} \cap U_{j} \neq \emptyset$. Assume that $\mathbf{U}$ is connected. Then $\Size{\cup^{k}_{i=1} U_{i}} \le \sum^{k}_{i=1}{\Size{U_{i}}}$. 
  \end{claim}
  \begin{proof}
    It suffices to prove that if $R_{1} \cap R_{2} \neq \emptyset$, then $\Size{R_{1} \cup R_{2}} \le \Size{R_{1}} + \Size{R_{2}}$. Let $a_{i,j} = \max_{v \in R_{j}}L_{i}(v)$. Then for any $w \in R_{1}\cap R_{2}$ we have:  
    \begin{equation*}
      \begin{split}
        \Size{R_{1} \cup R_{2}} &= \sum_{i}^{d+1}{\max{ a_{i,1},a_{i,2} } } =  \sum_{i}^{d+1}{a_{i,1} + a_{i,2} - \min{ a_{i,1},a_{i,2} } } \\
        & \le   \sum_{i}^{d+1}{a_{i,1} + a_{i,2} - L_{i}(w) } = \Size{R_{1}} + \Size{R_{2}} - \Size{w}  \\
        &  = \Size{R_{1}} + \Size{R_{2}} 
      \end{split}
    \end{equation*}
  \end{proof}

\begin{definition}
  We say that a Pauli operator $P$ is an \textit{$X$-type error} if it consists of a multiplication of single-qubit Pauli $X$ and identities. Similarly, we say that $P$ is a \textit{$Z$-type error} if it consists of a multiplication of single-qubit Pauli $Z$ and identities. For a binary assignment over the $d^{\prime}$-faces, $z : \mathcal{F}^{(i)} \rightarrow \mathbb{F}_{2}$ we say that $z$ is the \textit{underline assignment} of $P$ if $z$ is the support of the non-trivial coordinates of $P$. We extend the size function to $X/Z$-types errors as follow: 
  \begin{enumerate}
    \item If $P$ is a $X$-type error, then $\Size{P} \colonequals \Size{z}$. 
    \item Else, namely $P$ is a $Z$-type error, then $\Size{P} \colonequals \Size{\phi(z)}$. 
  \end{enumerate}
\end{definition}


\begin{claim}
  The following hold:
  \begin{enumerate}
    \item $ \phi \circ  H^{\otimes n} $ sends $X$-type errors to the same Size $Z$-type errors, and vice versa.  
    \item  $\mathbf{CX}$ and $\varphi$ send $X$-type errors to the same size $X$-type errors. Similarly, they send $Z$-types errors to the same size $Z$-type errors.  
  \end{enumerate}
\end{claim}
\begin{proof}
  Let $z$ be a binary assignment over the $d^{\prime}$-faces, $z : \mathcal{F}^{(i)} \rightarrow \mathbb{F}_{2}$, and let $P$ be the $X$-type error which $z$ is its underline assignment. Let $ Q = \phi \circ H^{\otimes n} P $, and denote by $z_{Q}$ the underline assignment of $Q$. 
\begin{equation*}
  \begin{split}
    Q \ket{\psi} & = \phi \circ H^{\otimes n} P \ket{\psi} = \phi \circ H^{\otimes n} X_{z} \ket{\psi} \\ 
    & = Z_{\phi (z)} \phi \circ H^{\otimes n} \ket{\psi} = Z_{\phi (z) } \ket{\psi_{2}} 
  \end{split}
\end{equation*}
Where $\ket{\psi_{2}}$ is a codeword. Thus, $Q$ is a $Z$-type error with underline assignment $z_{Q} = \phi(z)$, thus: 
\begin{equation*}
  \begin{split}
    \Size{Q} = \Size{\phi(z_{Q})} = \Size{\phi^{2}(z) } = \Size{z} = \Size{P}
  \end{split}
\end{equation*}


For $\mathbf{CX}$, we consider two registers, $R_1$ and $R_2$. The $\mathbf{CX}$ acts transversally on the registers, and it's clear, by commutation relations, that it duplicates $X$-type errors from $R_1$ to their counterparts on $R_2$. Similarly, $\mathbf{CX}$ duplicates $Z$-type errors from $R_2$ into their counterparts on $R_1$.


Lastly, for $\varphi$, denote the permutation that induces $\varphi$ by $\pi : S \rightarrow S$. Observes that for any subsets of vertices $U\subset \mathcal{G}$, we have:   
\begin{equation*}
  \begin{split}
    \sum_{g \in S}L_{g}(\varphi(U) )  &= \sum_{g \in S}L_{\pi(g)}(U ) = \sum_{g \in \pi(S)}L_{g}(U ) =\sum_{g \in S}L_{g}(U ) \\
    L_{d+1}(\varphi(U) ) &= \max_{v \in \varphi(U) } - \sum_{g \in S} v_{g} = \max_{v \in U } - \sum_{g \in S} v_{\pi(g)} = L_{d+1}(U)
  \end{split}
\end{equation*}
Namely, the size function is invariant to $\varphi$-rotations, $\Size{\varphi(U)} = \Size{U}$. Hence:
\begin{equation*}
  \begin{split}
    \Size{ \varphi (P) } &= \Size{ \varphi (X_{z}) } = \Size{ X_{\varphi(z)} } = \Size{ X_{z} } = \Size{ P }
  \end{split}
\end{equation*}
\end{proof}



\begin{claim}
  Let $R_{1}, R_{2}$ be two registers, then any logical operator in $\mathcal{O}$ sends error cluster to an error cluster as follows: 
\end{claim}



\section{Toric Encoded Circuits}

    \begin{definition}[Troic Encoded Circuit]
      Let $C$ be a quantum circuit. We say that $C$ is \textit{Toric encoded} if the following holds:
      \begin{enumerate}
          \item There is a Toric code $\mathcal{C}$ such that the qubits of $\mathcal{C}$ can be partitioned into registers $R_{1}, \ldots, R_{k}$, and each of the registers is encoded by $\mathcal{C}$.   
          \item The gates in $C$ with even time coordinates are sweep decoding gates.
        \end{enumerate} 
        For a set of faults $\mathcal{E}$ the gates in $C[\mathcal{E}]$ can be decomposed to cycles, of the following form:
        \begin{equation*}
          \begin{split}
            C =     D_{l}P_{l}\Lambda_{l},..,D_{1} P_{1} \Lambda_{1}
          \end{split}
        \end{equation*}
    \end{definition}

  \begin{definition}[Base graph $G_{o}$]
    Let $C$ be a quantum circuit, at depth $d$, with registers $R_{1}, \ldots, R_{k}$, all of which are encoded by Toric codes with the same parameters (i.e., dimension, length, etc.). Let's denote by $\mathcal{G}_{1}, .., \mathcal{G}_{k}$ the Toric complexes that induce the codes. We define the base graph of $C$, denoted by $G_{o} = (V_{o},E_{o})$ as follows: 
    \begin{enumerate}
      \item The vertices $V_{o}$ are tuples, where the first coordinates correspond to a vertex in the union over the vertices in $\mathcal{G}_{1}, \ldots, \mathcal{G}_{k}$, and the second coordinate is associated with a time. Namely:
        \begin{equation*}
          \begin{split}
            V_{o} \colonequals \bigsqcup_{i} V(\mathcal{G}_{i}) \times [d]
          \end{split}
        \end{equation*}
      \item The edges $E_{o}$ are derived from the original edges in $\mathcal{G}_{1}, \ldots, \mathcal{G}_{k}$, with the addition that any two vertices in preceding times $t$ and $t+1$ that support qubits (faces) such that $C$ acts non-trivially on those qubits at time $t$ are also connected by an edge. Namely:
        \begin{equation*}
          \begin{split}
            E_{o} & \colonequals \big\{ \{(v,t), (u,t^{\prime}) \} :  \text{ if }     \\ 
              & \text{either there exists } i \text{ such }  \{v,u\} \in \mathcal{G}_{i} \text{ and } t^{\prime} = t,t+1 \\
             & \text{or there exists } u_i = (g_{i},l_{i}) \in C \text{ such }  v,u \in l_{i,2} \text{ and } t = l_{i,1}, t^{\prime} = l_{i,1}+1  \big\} 
          \end{split}
        \end{equation*}
    \end{enumerate} 
\ctt{Here we should add that a vertex $v$ belongs to location $l$ if there is a qubit $f$ that $v$ belongs to its associated face. }

  \end{definition}




  \begin{definition}[Time Graph of $C(\mathcal{E})$]
    For a quantum circuit $C$ and set of faults $\mathcal{E}$, let us denote the base graph of $C[\mathcal{E}]$ by $G_{o}$. For each time $t$, denote the deviation induced by $\mathcal{E}$ at time $t$ by $\mathcal{D}_{t}$, so $\partial\mathcal{D}_{t}$ is the syndrome that the qubits in the deviation admit. We define the history graph $\mathcal{H} = ( V_{\mathcal{H}}, E_{\mathcal{H}})$ as follows: 
    \begin{enumerate}
      \item The vertices $V_{\mathcal{H}}$ are taken to be the vertices of $G_{o}$, namely tuples $(v,t) \in V_{o}$, that satisfy that $v$ belongs to a face associated with a check, which, at time $t$, belongs to the $\partial \mathcal{D}_{t}$. Namely:
        \begin{equation*}
          \begin{split}
            \big\{ v_{t} : v_{t}=(v,t) \in V_{o} \text{ and there exists a check } f \in C \text{ such } v \in f \in \partial \mathcal{D}_{t} \big\}
          \end{split}
        \end{equation*}
      \item The edges $E_{\mathcal{H}}$ is partitioned into two types $A_{\mathcal{H}}$ and $F_{\mathcal{H}}$, where $A_{\mathcal{H}}$ are the edges in $E_{o}$ restricted to $V_{\mathcal{H}}$ and $F_{\mathcal{H}}$ are defined to be the . Namely:
        \begin{equation*}
          \begin{split}
            A_{\mathcal{H}} &= \big\{ e = \{u,v\} \in E_{o} :  u,v \in V_{\mathcal{H}} \text{ and } \Time{v} = \Time{u}\pm 1 \big\} \\
            F_{\mathcal{H}} &= \big\{ e = \{u,v\} \in E_{o} : \text{ exists } w \in V_{\mathcal{H}} \text{ such } \{v,w\},\{u,w\} \in A_{\mathcal{H}} \big\} 
          \end{split}
        \end{equation*}
    \end{enumerate}
    By definition of the construction, $G_{\mathcal{H}}$ is a time-graph.
\end{definition}



  \begin{definition}[Size Invariant Action]
\ctt{Here we classify both transversal gates, and rotation on the Toric. }
An action $g$ is be said size invariant action if, for any assignment $z$, it holds that $\Size{gz} = \Size{z}$. 
  \end{definition}

  \begin{claim}
    The sweep rule is Size $\xi$-shrinking.
  \end{claim}
  \begin{proof}
    
  \end{proof}

  \ctt{Eventually, we might not used the projected graph, and then we will split to history of the slices lemma for 3 parts. The first is when time $=3t$, in that we apply an invariant action so. }

